\section{Заключение}

В ходе выполнения лабораторной работы с помощью метода наименьших квадратов были оценены коэффициенты регрессии вида $f(x) = w_0 + w_1 x + w_2 x^2$. При помощи выборки объёма \( n=20 \) были определы:

\begin{itemize}
    \item вектор неизвестных параметров:
    \[
        \hat{\vec{\omega}} = \begin{pmatrix} 4.94827296 & -3.93485988 & 0.98541809 \end{pmatrix}^{T}
    \]
    \item значение среднеквадратического отклонения:
    \[
        \hat{\sigma} \approx 0.0803965
    \]
\end{itemize}

Визуализация результатов на графике показала, что построенная кривая регрессии качественно аппроксимирует исходные точки выборки $(x, t)$.

Для автоматизации матричных вычислений \textit{(формирования матрицы плана $\Phi$, расчёта оценки МНК и вычисления СКО)} использовались скрипты, написанные на языке \textit{Python} с использованием библиотеки \textit{NumPy}. Их листинги приведены в теле отчёта.